\documentclass[10pt,a4paper]{amsart}
\usepackage{amsmath,amssymb,amsthm,graphicx}
% Retained preamble for the Stacks Project chapter snapshot in this directory.
% The chapter's own preamble.tex is a copy of the shared Stacks Project preamble; the block
% below is sliced MECHANICALLY, line for line, from it. Nothing was added, removed or
% redefined; the wrapper class, the xr-hyper/hyperref wiring and the externaldocument list
% are omitted because the assembler supplies the document class.
%
% lines taken from preamble.tex: 3, 5, 9, 11, 12, 156, 157, 158, ... (51 lines total)
%
% package coverage restored from the chapter's own preamble (2026-10-13): xy, xy-2cell

\usepackage[all]{xy}
\xyoption{2cell}
\UseAllTwocells

\newenvironment{reference}{\comment}{\endcomment}
\newenvironment{slogan}{\comment}{\endcomment}
\newenvironment{history}{\comment}{\endcomment}
\theoremstyle{plain}
\newtheorem{theorem}[subsection]{Theorem}
\newtheorem{proposition}[subsection]{Proposition}
\newtheorem{lemma}[subsection]{Lemma}
\theoremstyle{definition}
\newtheorem{definition}[subsection]{Definition}
\newtheorem{example}[subsection]{Example}
\newtheorem{exercise}[subsection]{Exercise}
\newtheorem{situation}[subsection]{Situation}
\theoremstyle{remark}
\newtheorem{remark}[subsection]{Remark}
\newtheorem{remarks}[subsection]{Remarks}
\numberwithin{equation}{subsection}
\def\lim{\mathop{\mathrm{lim}}\nolimits}
\def\colim{\mathop{\mathrm{colim}}\nolimits}
\def\Spec{\mathop{\mathrm{Spec}}}
\def\Hom{\mathop{\mathrm{Hom}}\nolimits}
\def\Ext{\mathop{\mathrm{Ext}}\nolimits}
\def\SheafHom{\mathop{\mathcal{H}\!\mathit{om}}\nolimits}
\def\SheafExt{\mathop{\mathcal{E}\!\mathit{xt}}\nolimits}
\def\Sch{\mathit{Sch}}
\def\Mor{\mathop{\mathrm{Mor}}\nolimits}
\def\Ob{\mathop{\mathrm{Ob}}\nolimits}
\def\Sh{\mathop{\mathit{Sh}}\nolimits}
\def\NL{\mathop{N\!L}\nolimits}
\def\CH{\mathop{\mathrm{CH}}\nolimits}
\def\proetale{{pro\text{-}\acute{e}tale}}
\def\etale{{\acute{e}tale}}
\def\QCoh{\mathit{QCoh}}
\def\Ker{\mathop{\mathrm{Ker}}}
\def\Im{\mathop{\mathrm{Im}}}
\def\Coker{\mathop{\mathrm{Coker}}}
\def\Coim{\mathop{\mathrm{Coim}}}
\DeclareMathSymbol{\boxtimes}{\mathbin}{AMSa}{"02}
\def\QCohstack{\mathcal{QC}\!\mathit{oh}}
\def\Cohstack{\mathcal{C}\!\mathit{oh}}
\def\Spacesstack{\mathcal{S}\!\mathit{paces}}
\def\Quotfunctor{\mathrm{Quot}}
\def\Hilbfunctor{\mathrm{Hilb}}
\def\Curvesstack{\mathcal{C}\!\mathit{urves}}
\def\Polarizedstack{\mathcal{P}\!\mathit{olarized}}
\def\Complexesstack{\mathcal{C}\!\mathit{omplexes}}
\def\Pic{\mathop{\mathrm{Pic}}\nolimits}
\def\Picardstack{\mathcal{P}\!\mathit{ic}}
\def\Picardfunctor{\mathrm{Pic}}
\def\Deformationcategory{\mathcal{D}\!\mathit{ef}}
\usepackage{cleveref}
\crefname{originalproposition}{original proposition}{original propositions}
\Crefname{originalproposition}{Original proposition}{Original propositions}
\makeatletter
\newlabel{categories-example-functor-groupoids}{{0049}{source}{}{sourcepointer-1}{}}
\newlabel{categories-example-functor-groupoids@cref}{{[originalproposition][0][]0049}{[1][0][]source}{}{}{}}
\makeatother
\crefname{originalproposition}{original proposition}{original propositions}
\Crefname{originalproposition}{Original proposition}{Original propositions}
\makeatletter
\newlabel{situation-quotient-stack}{{04WL}{source}{}{sourcepointer-2}{}}
\newlabel{situation-quotient-stack@cref}{{[originalproposition][0][]04WL}{[1][0][]source}{}{}{}}
\makeatother
\crefname{originalproposition}{original proposition}{original propositions}
\Crefname{originalproposition}{Original proposition}{Original propositions}
\makeatletter
\newlabel{spaces-groupoids-definition-quotient-stack}{{044Q}{source}{}{sourcepointer-3}{}}
\newlabel{spaces-groupoids-definition-quotient-stack@cref}{{[originalproposition][0][]044Q}{[1][0][]source}{}{}{}}
\makeatother
\crefname{originalproposition}{original proposition}{original propositions}
\Crefname{originalproposition}{Original proposition}{Original propositions}
\makeatletter
\newlabel{spaces-groupoids-lemma-quotient-stack-morphisms}{{044V}{source}{}{sourcepointer-4}{}}
\newlabel{spaces-groupoids-lemma-quotient-stack-morphisms@cref}{{[originalproposition][0][]044V}{[1][0][]source}{}{}{}}
\makeatother
\crefname{originalproposition}{original proposition}{original propositions}
\Crefname{originalproposition}{Original proposition}{Original propositions}
\makeatletter
\newlabel{spaces-groupoids-section-explicit-quotient-stacks}{{04MB}{source}{}{sourcepointer-5}{}}
\newlabel{spaces-groupoids-section-explicit-quotient-stacks@cref}{{[originalproposition][0][]04MB}{[1][0][]source}{}{}{}}
\makeatother
\crefname{originalproposition}{original proposition}{original propositions}
\Crefname{originalproposition}{Original proposition}{Original propositions}
\makeatletter
\newlabel{stacks-lemma-characterize-essentially-surjective-when-ff}{{046N}{source}{}{sourcepointer-6}{}}
\newlabel{stacks-lemma-characterize-essentially-surjective-when-ff@cref}{{[originalproposition][0][]046N}{[1][0][]source}{}{}{}}
\makeatother
\crefname{originalproposition}{original proposition}{original propositions}
\Crefname{originalproposition}{Original proposition}{Original propositions}
\makeatletter
\newlabel{stacks-lemma-stackify-groupoids-universal-property}{{0436}{source}{}{sourcepointer-7}{}}
\newlabel{stacks-lemma-stackify-groupoids-universal-property@cref}{{[originalproposition][0][]0436}{[1][0][]source}{}{}{}}
\makeatother
\title{Canonical equivalence of the quotient stack and the torsor quotient stack}
\author{The Stacks Project Authors}
\date{stacks-project:examples-stacks.tex:local-snapshot-20261012; GNU FDL 1.2 or later}
\begin{document}
\maketitle
\noindent\textit{Editorial note.} Statement span is lines 1527-1536 of sources/stacks-examples-stacks/pinned-original-source.tex; proof span is lines 1538-1673, the complete author proof (134 body lines).\par\medskip
\noindent\textit{Editorial note.} Navigation stubs: 7 cited-but-not-included labels are supplied as non-load-bearing pointers with their pinned tags where the snapshot tag index has them.\par\medskip
\noindent\textit{Editorial note.} extra\_packages is empty on purpose: the retained preamble already loads xy with options and a second declaration clashes in the assembler wrapper.\par\medskip
\noindent\textit{Editorial note.} Text is retained verbatim, including the author's own typos and placeholders.\par\medskip
\noindent\textit{Editorial note.} Queue index 0 of queue\_examples-stacks.json (pb=4979, stmt 1527-1536, proof 1538-1673, halfA).\par\medskip
\noindent\textit{Editorial note.} stacks-examples-stacks wave-1 (exs1): fresh adjudication accept (2026-10-14); pool per Chief 09:4x authorization (idle-line throughput over kun's backup note, SMC precedent); examples-stacks.tex 27 labels/20 with proof judged 9 accept / 11 reject; claim\_key v2 stacks:examples-stacks.tex:<label>.\par\medskip
\noindent\textit{Allowed standard prerequisites (editorial).}\begin{itemize}
\item Explicit prerequisite boundary. The cited inputs are prerequisites; the argument in the retained span is not replaced by any cited result.
\end{itemize}\medskip
\ifdefined\hypertarget\hypertarget{sourcepointer-1}{}\fi
\noindent\textit{Source reference.} Original proposition 0049: sources/stacks-examples-stacks/pinned-original-source.tex, source lines 1673--1673. Cited by the retained statement or proof but not part of the collected spans; supplied only as a navigation pointer. The argument does not depend on its content.\par\medskip
\ifdefined\hypertarget\hypertarget{sourcepointer-2}{}\fi
\noindent\textit{Source reference.} Original proposition 04WL: sources/stacks-examples-stacks/pinned-original-source.tex, source lines 1418--1426. Cited by the retained statement or proof but not part of the collected spans; supplied only as a navigation pointer. The argument does not depend on its content.\par\medskip
\ifdefined\hypertarget\hypertarget{sourcepointer-3}{}\fi
\noindent\textit{Source reference.} Original proposition 044Q: sources/stacks-examples-stacks/pinned-original-source.tex, source lines 1673--1673. Cited by the retained statement or proof but not part of the collected spans; supplied only as a navigation pointer. The argument does not depend on its content.\par\medskip
\ifdefined\hypertarget\hypertarget{sourcepointer-4}{}\fi
\noindent\textit{Source reference.} Original proposition 044V: sources/stacks-examples-stacks/pinned-original-source.tex, source lines 1673--1673. Cited by the retained statement or proof but not part of the collected spans; supplied only as a navigation pointer. The argument does not depend on its content.\par\medskip
\ifdefined\hypertarget\hypertarget{sourcepointer-5}{}\fi
\noindent\textit{Source reference.} Original proposition 04MB: sources/stacks-examples-stacks/pinned-original-source.tex, source lines 1673--1673. Cited by the retained statement or proof but not part of the collected spans; supplied only as a navigation pointer. The argument does not depend on its content.\par\medskip
\ifdefined\hypertarget\hypertarget{sourcepointer-6}{}\fi
\noindent\textit{Source reference.} Original proposition 046N: sources/stacks-examples-stacks/pinned-original-source.tex, source lines 1673--1673. Cited by the retained statement or proof but not part of the collected spans; supplied only as a navigation pointer. The argument does not depend on its content.\par\medskip
\ifdefined\hypertarget\hypertarget{sourcepointer-7}{}\fi
\noindent\textit{Source reference.} Original proposition 0436: sources/stacks-examples-stacks/pinned-original-source.tex, source lines 1673--1673. Cited by the retained statement or proof but not part of the collected spans; supplied only as a navigation pointer. The argument does not depend on its content.\par\medskip
\begin{proposition}
\label{proposition-equal-quotient-stacks}
In
Situation \ref{situation-quotient-stack}
there exists a canonical equivalence
$$
[X/G] \longrightarrow [[X/G]]
$$
of stacks in groupoids over $(\Sch/S)_{fppf}$.
\end{proposition}

\begin{proof}
We write this out in detail, to make sure that all the definitions
work out in exactly the correct manner.
Recall that $[X/G]$ is the quotient stack
associated to the groupoid in algebraic spaces
$(X, G \times_B X, s, t, c)$, see
Groupoids in Spaces,
Definition \ref{spaces-groupoids-definition-quotient-stack}.
This means that $[X/G]$ is the stackification of the
category fibred in groupoids $[X/_{\!p}G]$ associated to the functor
$$
(\Sch/S)_{fppf} \longrightarrow \textit{Groupoids},
\quad
U \longmapsto (X(U), G(U) \times_{B(U)} X(U), s, t, c)
$$
where $s(g, x) = x$, $t(g, x) = a(g, x)$, and
$c((g, x), (g', x')) = (m(g, g'), x')$. By the construction of
Categories, Example \ref{categories-example-functor-groupoids}
an object of $[X/_{\!p}G]$ is a pair $(U, x)$ with $x \in X(U)$
and a morphism $(f, g) : (U, x) \to (U', x')$ of $[X/_{\!p}G]$
is given by a morphism of schemes $f : U \to U'$ and an element
$g \in G(U)$ such that $a(g, x) = x' \circ f$.
Hence we can define a $1$-morphism of stacks in groupoids
$$
F_p : [X/_{\!p}G] \longrightarrow [[X/G]]
$$
by the following rules: On objects we set
$$
F_p(U, x) =
(U, \pi \circ x, G \times_{B, \pi \circ x} U, a \circ (\text{id}_G \times x))
$$
This makes sense because the diagram
$$
\xymatrix{
G \times_{B, \pi \circ x} U \ar[d] \ar[r]_{\text{id}_G \times x} &
G \times_{B, \pi} X \ar[r]_-a &
X \ar[d]^\pi \\
U \ar[rr]^{\pi \circ x} & & B
}
$$
commutes, and the two horizontal arrows are $G$-equivariant if we think
of the fibre products as trivial $G$-torsors over $U$, resp.\ $X$.
On morphisms $(f, g) : (U, x) \to (U', x')$ we set
$F_p(f, g) = (f, R_{g^{-1}})$ where $R_{g^{-1}}$ denotes right translation
by the inverse of $g$. More precisely, the morphism
$F_p(f, g) : F_p(U, x) \to F_p(U', x')$ is given by the cartesian diagram
$$
\xymatrix{
G \times_{B, \pi \circ x} U \ar[d] \ar[r]_{R_{g^{-1}}} &
G \times_{B, \pi \circ x'} U' \ar[d] \\
U \ar[r]^f & U'
}
$$
where $R_{g^{-1}}$ on $T$-valued points is given by
$$
R_{g^{-1}}(g', u) = (m(g', i(g(u))), f(u))
$$
To see that this works we have to verify that
$$
a \circ (\text{id}_G \times x)
=
a \circ (\text{id}_G \times x') \circ R_{g^{-1}}
$$
which is true because the right hand side applied to the $T$-valued point
$(g', u)$ gives the desired equality
\begin{align*}
a((\text{id}_G \times x')(m(g', i(g(u))), f(u)))
& =
a(m(g', i(g(u))), x'(f(u))) \\
& =
a(g', a(i(g(u)), x'(f(u)))) \\
& =
a(g', x(u))
\end{align*}
because $a(g, x) = x' \circ f$ and hence $a(i(g), x' \circ f) = x$.

\medskip\noindent
By the universal property of stackification from
Stacks, Lemma \ref{stacks-lemma-stackify-groupoids-universal-property}
we obtain a canonical extension $F : [X/G] \to [[X/G]]$ of the $1$-morphism
$F_p$ above. We first prove that $F$ is fully faithful.
To do this, since both source and target are stacks in groupoids,
it suffices to prove that the $\mathit{Isom}$-sheaves are identified
under $F$. Pick a scheme $U$ and objects $\xi, \xi'$ of
$[X/G]$ over $U$. We want to show that
$$
F :
\mathit{Isom}_{[X/G]}(\xi, \xi')
\longrightarrow
\mathit{Isom}_{[[X/G]]}(F(\xi), F(\xi'))
$$
is an isomorphism of sheaves. To do this it suffices to work locally
on $U$, and hence we may assume that $\xi, \xi'$ come from objects
$(U, x)$, $(U, x')$ of $[X/_{\!p}G]$ over $U$; this follows directly
from the construction of the stackification, and it is also worked
out in detail in
Groupoids in Spaces,
Section \ref{spaces-groupoids-section-explicit-quotient-stacks}.
Either by directly using the description of morphisms in
$[X/_{\!p}G]$ above, or using
Groupoids in Spaces,
Lemma \ref{spaces-groupoids-lemma-quotient-stack-morphisms}
we see that in this case
$$
\mathit{Isom}_{[X/G]}(\xi, \xi') =
U \times_{(x, x'), X \times_S X, (s, t)} (G \times_B X)
$$
A $T$-valued point of this fibre product corresponds to a pair
$(u, g)$ with $u \in U(T)$, and $g \in G(T)$ such that
$a(g, x \circ u) = x' \circ u$. (Note that this implies
$\pi \circ x \circ u = \pi \circ x' \circ u$.)
On the other hand, a $T$-valued
point of $\mathit{Isom}_{[[X/G]]}(F(\xi), F(\xi'))$ by definition
corresponds to a morphism $u : T \to U$ such that
$\pi \circ x \circ u = \pi \circ x' \circ u : T \to B$ and an isomorphism
$$
R :
G \times_{B, \pi \circ x \circ u} T
\longrightarrow
G \times_{B, \pi \circ x' \circ u} T
$$
of trivial $G_T$-torsors compatible with the given maps to $X$.
Since the torsors are trivial we see that $R = R_{g^{-1}}$
(right multiplication) by some $g \in G(T)$. Compatibility with the maps
$a \circ (1_G, x \circ u), a \circ (1_G, x' \circ u) : G \times_B T \to X$
is equivalent to the condition that $a(g, x \circ u) = x' \circ u$.
Hence we obtain the desired equality of $\mathit{Isom}$-sheaves.

\medskip\noindent
Now that we know that $F$ is fully faithful we see that
Stacks, Lemma \ref{stacks-lemma-characterize-essentially-surjective-when-ff}
applies. Thus to show that $F$ is an equivalence it suffices
to show that objects of $[[X/G]]$ are fppf locally in the essential image
of $F$. This is clear as fppf torsors are fppf locally trivial, and hence
we win.
\end{proof}

\end{document}
